% !TEX program = xelatex
\documentclass[UTF8,12pt,a4paper]{ctexart}

\usepackage[a4paper,margin=2.5cm]{geometry}
\usepackage{amsmath,amssymb,bm}
\usepackage{booktabs}
\usepackage{graphicx}
\usepackage{float}
\usepackage{subcaption}
\usepackage{xcolor}
\usepackage{hyperref}
\usepackage{enumitem}

\graphicspath{{figures/}}

\hypersetup{
    colorlinks=true,
    linkcolor=blue!55!black,
    citecolor=blue!55!black,
    urlcolor=blue!55!black
}

\setlist[itemize]{leftmargin=2em}

\newcommand{\caseplots}[3]{
\begin{figure}[H]
    \centering
    \begin{subfigure}{0.48\textwidth}
        \centering
        \includegraphics[width=\linewidth]{#1_rho.png}
        \caption{密度 $\rho$}
    \end{subfigure}
    \hfill
    \begin{subfigure}{0.48\textwidth}
        \centering
        \includegraphics[width=\linewidth]{#1_u.png}
        \caption{速度 $u$}
    \end{subfigure}
    \vspace{0.6em}

    \begin{subfigure}{0.62\textwidth}
        \centering
        \includegraphics[width=\linewidth]{#1_p.png}
        \caption{压力 $p$}
    \end{subfigure}
    \caption{#2：MacCormack 数值解与精确解对比。#3}
    \label{fig:#1}
\end{figure}
}

\title{一维 Riemann 问题的 MacCormack 数值结果与讨论}
\author{邵桂博、杨振宇}
\date{\today}

\begin{document}

\maketitle

\begin{abstract}
本文用 MacCormack 格式计算一维 Euler 方程的七组 Riemann 问题，并与精确 Riemann 解进行对比。分别考察人工粘性的有无、经验系数 $\beta$、CFL 数、开关变量以及网格加密对结果的影响。计算结果表明，MacCormack 格式在光滑区域和接触间断问题上可以得到较好的结果，但在激波、强膨胀和近真空区域中会出现明显的色散振荡，甚至产生非物理状态。人工粘性可以提高间断附近的稳定性，但过大的人工粘性会抹宽波系，因此它不是单纯提高精度的手段，而是在稳定性和分辨率之间作折中。

\textbf{关键词：} Riemann 问题；MacCormack 格式；人工粘性；Euler 方程；
\end{abstract}

\tableofcontents
\newpage

\section{定义说明}

在正式写格式之前，先说明本文后面反复出现的几个参数和概念。
\begin{itemize}
    \item \textbf{人工粘性}：这里不是气体真实的物理粘性，而是人为加到数值格式里的耗散项。它的作用是压制激波、接触间断附近的非物理振荡，但加得太强会把间断抹宽。
    \item \textbf{sensor}：本文也称为开关变量，意思是用哪个物理量来判断哪里需要加人工粘性。程序中可以选 $\rho$、$u$ 或 $p$。例如 sensor=$\rho$ 表示用密度变化计算开关函数，sensor=$p$ 表示用压力变化计算开关函数。
    \item \textbf{$\beta$}：人工粘性的经验系数。$\beta$ 越大，人工粘性通常越强，稳定性可能更好，但数值扩散也会更明显。
    \item \textbf{CFL 数}：控制时间步长的无量纲参数。本文中时间步长按局部最大波速调整，CFL 越大，每一步走得越大，但稳定性风险也会增加。
    \item \textbf{$N_x$}：一维网格点数。$N_x$ 越大，空间分辨率越高，但计算量也越大。
    \item \textbf{mode}：MacCormack 的差分方向选择。mode 1 表示预估步 FTBS、校正步 FTFS；mode 2 表示预估步 FTFS、校正步 FTBS；mode 3 和 mode 4 表示两种方向交替使用。
    \item \textbf{completed/failed}：批量计算中的状态记录。completed 表示该参数组合能算到终止时间；failed 表示计算中出现负密度、负压力或 NaN 等非物理状态。
\end{itemize}

本文使用的误差定义为
\begin{align}
    L_1(\phi)
    &=
    \sum_{i=0}^{N-2}
    \frac{|e_i|+|e_{i+1}|}{2}
    (x_{i+1}-x_i),
    \\
    L_2(\phi)
    &=
    \left[
    \sum_{i=0}^{N-2}
    \frac{e_i^2+e_{i+1}^2}{2}
    (x_{i+1}-x_i)
    \right]^{1/2},
    \\
    L_\infty(\phi)
    &=
    \max_i |e_i|.
\end{align}
其中 $L_1$ 反映整体偏差，$L_2$ 对较大的局部误差更敏感，$L_\infty$ 是最大点误差。对于含间断的问题，$L_\infty$ 对激波和接触间断的位置误差非常敏感，所以本文不只用最大误差判断格式优劣。

\section{算例与批量实验设置}

七组 Riemann 问题均在区间 $[0,1]$ 上计算，初始间断位置为 $x_0=0.5$。基准网格为 $N_x=501$。除 LAX 算例取 $t=0.16$、第三和第四组取 $t=0.15$ 外，其余算例取 $t=0.2$。

\begin{table}[H]
    \centering
    \caption{七组算例的基准参数}
    \begin{tabular}{clccccc}
        \toprule
        Case & 名称 & $t$ & $N_x$ & CFL & sensor & $\beta$ \\
        \midrule
        1 & Sod 激波管 & 0.20 & 501 & 0.5 & $\rho$ & 0.5 \\
        2 & LAX 问题 & 0.16 & 501 & 0.5 & $\rho$ & 0.5 \\
        3 & 亚声速双膨胀 & 0.15 & 501 & 0.2 & $\rho$ & 0.5 \\
        4 & Sjogreen 超声速膨胀 & 0.15 & 501 & 0.2 & $u$ & 1.0 \\
        5 & 接触间断双膨胀 & 0.20 & 501 & 0.5 & $\rho$ & 0.5 \\
        6 & 双激波间断 & 0.20 & 501 & 0.5 & $\rho$ & 0.5 \\
        7 & 纯接触间断 & 0.20 & 501 & 0.5 & $\rho$ & 0.5 \\
        \bottomrule
    \end{tabular}
\end{table}

为了讨论人工粘性和参数影响，每个算例还做了批量实验。这部分不是手工一组一组输入完成的，而是用
\verb|postprocess/run_report_matrix.py|
自动调用 C 程序。脚本每次改变一组参数，运行 MacCormack 程序，再把输出文件和误差表保存下来。
主要变化量为：
\begin{equation}
    \beta=0.25,0.5,0.75,1.0,1.5,
    \qquad
    \mathrm{CFL}=0.2,0.5,0.8,
\end{equation}
开关变量取 $\rho,u,p$，网格数取 $101,201,501,1001$，并额外加入关闭人工粘性的情况。组数的来源为
\begin{equation}
    N_{\mathrm{on}}
    =
    5\times3\times3\times4
    =
    180 ,
    \qquad
    N_{\mathrm{off}}
    =
    3\times4
    =
    12 .
\end{equation}
其中 $N_{\mathrm{on}}$ 是开启人工粘性时的参数组合数，分别对应 $\beta$、sensor、CFL 和网格数；$N_{\mathrm{off}}$ 是关闭人工粘性时的参数组合数，此时 $\beta$ 和 sensor 不再起作用，所以只改变 CFL 和网格数。因此每个算例共统计
\begin{equation}
    N_{\mathrm{total}}=180+12=192
\end{equation}
组参数。运行结果见表 \ref{tab:status}。

\begin{table}[H]
    \centering
    \caption{七组批量实验的完成情况}
    \label{tab:status}
    \begin{tabular}{clrr}
        \toprule
        Case & 名称 & 完成组数 & 失败组数 \\
        \midrule
        1 & Sod 激波管 & 96 & 96 \\
        2 & LAX 问题 & 192 & 0 \\
        3 & 亚声速双膨胀 & 72 & 120 \\
        4 & Sjogreen 超声速膨胀 & 8 & 184 \\
        5 & 接触间断双膨胀 & 192 & 0 \\
        6 & 双激波间断 & 192 & 0 \\
        7 & 纯接触间断 & 192 & 0 \\
        \bottomrule
    \end{tabular}
\end{table}

这里的失败不是程序没有运行，而是计算过程中出现负密度、负压力或 NaN 等非物理状态。对于激波和强膨胀问题，这些失败本身也反映了格式的稳定性限制。

\section{基准结果对比}

\subsection{Case 1：Sod 激波管}

Sod 问题是最常用的激波管测试之一，波系中包含稀疏波、接触间断和激波。基准计算的误差见表 \ref{tab:baseline_error}。密度和压力的整体误差较小，速度的 $L_\infty$ 较大，主要来自激波和接触间断附近的局部偏差。

\caseplots{case1}{Case 1 Sod 激波管}{数值解基本能抓住三种波，但间断附近仍有被抹宽的现象。}

\subsection{Case 2：LAX 问题}

LAX 问题的波强比 Sod 更大，局部最大误差明显增加。与 Sod 不同，本组 192 个参数组合全部完成，即使关闭人工粘性也可以推进到终止时间。

\caseplots{case2}{Case 2 LAX 问题}{激波和接触间断附近的误差更明显，人工粘性会使间断变宽。}

\subsection{Case 3：亚声速双膨胀问题}

本算例中两侧速度相背，中心区域形成膨胀结构。直接使用 CFL=0.5 时容易失败，因此基准计算改为 CFL=0.2。这个结果说明膨胀问题虽然没有强激波，但仍可能因为局部状态变化导致 MacCormack 推进不稳定。

\caseplots{case3}{Case 3 亚声速双膨胀问题}{降低 CFL 后可以得到稳定结果，整体误差比强激波问题小。}

\subsection{Case 4：Sjogreen 超声速膨胀问题}

第四组是最困难的算例。超声速膨胀会在中间区域产生接近真空的状态，密度和压力很低，数值格式稍有振荡就可能导致负压或负密度。批量实验中只有 8 组完成，稳定组合基本集中在 velocity sensor、CFL=0.2、$\beta=1.0$ 或 $1.5$。

\caseplots{case4}{Case 4 Sjogreen 超声速膨胀问题}{近真空区域对格式非常敏感，能算通的参数范围明显变窄。}

\subsection{Case 5：接触间断双膨胀问题}

第五组比第四组温和，192 组参数全部完成。基准结果中密度误差主要集中在接触间断附近，速度和压力误差较小。适量人工粘性可以降低局部振荡，但过大时仍会增加扩散误差。

\caseplots{case5}{Case 5 接触间断双膨胀问题}{整体稳定性较好，误差主要由接触间断附近的密度变化贡献。}

\subsection{Case 6：双激波间断问题}

第六组包含两个激波结构。批量实验全部完成，说明在当前终止时间和初值下，程序能稳定处理该组强压缩问题。不过人工粘性增大后，密度 $L_1$ 误差总体上升，说明激波被进一步抹宽。

\caseplots{case6}{Case 6 双激波间断问题}{格式能够稳定捕捉双激波，但激波附近的局部误差仍然明显。}

\subsection{Case 7：纯接触间断问题}

纯接触间断中压力和速度理论上保持一致，主要变化量是密度。计算结果也符合这一点：速度和压力误差接近机器精度，密度误差占主导。这个算例说明不同格式有各自适合的问题，MacCormack 在纯接触问题上并不一定差。

\caseplots{case7}{Case 7 纯接触间断问题}{速度和压力几乎保持不变，误差主要来自密度间断的数值扩散。}

\begin{table}[H]
    \centering
    \caption{七组基准结果的误差}
    \label{tab:baseline_error}
    \begin{tabular}{clccc}
        \toprule
        Case & 名称 & $\rho$ 的 $L_1$ & $u$ 的 $L_1$ & $p$ 的 $L_1$ \\
        \midrule
        1 & Sod 激波管 & $4.117286\times10^{-3}$ & $1.160388\times10^{-2}$ & $3.580555\times10^{-3}$ \\
        2 & LAX 问题 & $2.643395\times10^{-2}$ & $3.426100\times10^{-2}$ & $4.336475\times10^{-2}$ \\
        3 & 亚声速双膨胀 & $1.656326\times10^{-3}$ & $6.557940\times10^{-3}$ & $6.834567\times10^{-3}$ \\
        4 & Sjogreen 超声速膨胀 & $5.393412\times10^{-3}$ & $1.690367\times10^{-2}$ & $2.531564\times10^{-3}$ \\
        5 & 接触间断双膨胀 & $3.277088\times10^{-3}$ & $2.346795\times10^{-3}$ & $1.269157\times10^{-3}$ \\
        6 & 双激波间断 & $1.050757\times10^{-2}$ & $6.922352\times10^{-3}$ & $8.577551\times10^{-3}$ \\
        7 & 纯接触间断 & $6.695606\times10^{-2}$ & $6.165206\times10^{-10}$ & $3.017485\times10^{-9}$ \\
        \bottomrule
    \end{tabular}
\end{table}

\section{MacCormack 方向模式对比}

MacCormack 格式本身有两个方向相反的写法。为了简单比较这件事，本文只选 Sod 算例做一组补充计算。选择 Sod 的原因是它同时包含稀疏波、接触间断和激波，比纯接触问题更能看出方向格式对间断附近误差的影响；同时它又不像第四组近真空问题那么敏感，便于保持其他参数不变。

本组计算固定
\[
    N_x=501,\qquad
    \mathrm{CFL}=0.5,\qquad
    \beta=0.5,\qquad
    \mathrm{sensor}=\rho,\qquad
    t=0.2 .
\]
只改变 MacCormack 的方向模式。mode 1 和 mode 2 表示每一步都使用同一种方向组合；mode 3 和 mode 4 表示两种方向交替使用。

\begin{table}[H]
    \centering
    \caption{Sod 算例中不同 MacCormack 方向模式的 $L_1$ 误差}
    \label{tab:mode_compare}
    \begin{tabular}{clccc}
        \toprule
        mode & 方向模式 & $\rho$ 的 $L_1$ & $u$ 的 $L_1$ & $p$ 的 $L_1$ \\
        \midrule
        1 & 固定 FTBS/FTFS & $4.117286\times10^{-3}$ & $1.160388\times10^{-2}$ & $3.580555\times10^{-3}$ \\
        2 & 固定 FTFS/FTBS & $3.495325\times10^{-3}$ & $8.858318\times10^{-3}$ & $2.895035\times10^{-3}$ \\
        3 & 交替，FTBS/FTFS 起步 & $3.768057\times10^{-3}$ & $1.014959\times10^{-2}$ & $3.200490\times10^{-3}$ \\
        4 & 交替，FTFS/FTBS 起步 & $3.732765\times10^{-3}$ & $1.010236\times10^{-2}$ & $3.169967\times10^{-3}$ \\
        \bottomrule
    \end{tabular}
\end{table}

从表 \ref{tab:mode_compare} 可以看到，在这组 Sod 参数下，mode 2 的三个变量 $L_1$ 都最小，两个交替方向的结果处在 mode 1 和 mode 2 之间，并且 mode 3 与 mode 4 的差别很小。这说明交替方向可以减弱单一方向带来的偏置，但并不保证误差一定比两个固定方向都小。后面的七组基准计算仍统一采用 mode 1，主要是为了让不同算例之间的比较更一致。

\section{人工粘性影响}

\subsection{有无人工粘性}

人工粘性的作用在不同算例中并不完全一样。Sod 问题中，关闭人工粘性后所有测试网格和 CFL 都失败，说明单纯加密网格不能消除 MacCormack 在激波附近的非物理振荡。LAX 问题则不同，关闭人工粘性仍然能稳定计算，并且在 $N_x=501,\ \mathrm{CFL}=0.5$ 时，关闭人工粘性的密度 $L_1$ 为 $1.897108\times10^{-2}$，小于基准人工粘性结果的 $2.643395\times10^{-2}$。

这个现象说明，人工粘性首先是稳定化手段，不是保证误差一定下降的手段。对容易振荡甚至失败的激波问题，它可以保证计算完成；对已经稳定的算例，它可能主要表现为额外扩散。

\subsection{$\beta$ 的影响}

表 \ref{tab:beta} 给出了固定基准网格和基准 CFL 时，不同 $\beta$ 对密度 $L_1$ 误差的影响。可以看到，$\beta$ 并不是越大越好。较小的 $\beta$ 有时能给出较低误差，但在 Sod 和强膨胀问题中，过小的 $\beta$ 可能不能稳定计算。

\begin{table}[H]
    \centering
    \caption{不同 $\beta$ 下的密度 $L_1$ 误差}
    \label{tab:beta}
    \begin{tabular}{clcccccc}
        \toprule
        Case & 名称 & $\beta=0$ & 0.25 & 0.5 & 0.75 & 1.0 & 1.5 \\
        \midrule
        1 & Sod & -- & failed & $4.1173$e-3 & $4.5992$e-3 & $4.9063$e-3 & $5.2862$e-3 \\
        2 & LAX & $1.8971$e-2 & $2.1129$e-2 & $2.6434$e-2 & $2.9019$e-2 & $3.0008$e-2 & $3.0698$e-2 \\
        3 & 亚声速双膨胀 & -- & failed & $1.6563$e-3 & $2.1734$e-3 & $2.6418$e-3 & $3.4366$e-3 \\
        4 & 超声速膨胀 & -- & failed & failed & failed & $5.3934$e-3 & $5.4027$e-3 \\
        5 & 接触双膨胀 & $5.0928$e-3 & $2.8642$e-3 & $3.2771$e-3 & $3.6870$e-3 & $4.0639$e-3 & $4.7022$e-3 \\
        6 & 双激波 & $8.2523$e-3 & $8.6533$e-3 & $1.0508$e-2 & $1.1293$e-2 & $1.1739$e-2 & $1.2456$e-2 \\
        7 & 纯接触 & $9.7251$e-2 & $6.2287$e-2 & $6.6956$e-2 & $7.2505$e-2 & $7.7483$e-2 & $8.6141$e-2 \\
        \bottomrule
    \end{tabular}
\end{table}

从表中可以总结出三点：
\begin{itemize}
    \item 对 Sod 和膨胀问题，$\beta$ 太小容易失稳；$\beta$ 足够大后可以计算，但误差通常随耗散增强而变大。
    \item 对 LAX 和双激波问题，关闭人工粘性或小 $\beta$ 的 $L_1$ 误差反而更低，说明过强人工粘性会降低间断分辨率。
    \item 对纯接触间断，$\beta=0.25$ 比无人工粘性更好，但继续增大 $\beta$ 后密度误差又增大。
\end{itemize}

\section{CFL、开关变量与网格影响}

\subsection{CFL 的影响}

人工粘性中含有 $\mathrm{CFL}(1-\mathrm{CFL})$。这个因子在 CFL=0.5 附近最大，但它不能单独决定稳定性。例如 Sod 问题中 CFL=0.2 和 CFL=0.8 的该因子相同，都是 0.16，但 CFL=0.2 可以完成，CFL=0.8 失败。这说明时间步长本身也会影响 MacCormack 推进的稳定性。

LAX 问题中固定 $\beta=0.5$、sensor=$\rho$、$N_x=501$ 时，CFL=0.8 的密度 $L_1$ 反而小于 CFL=0.2 和 0.5。这也说明 CFL 对误差的影响不是简单单调关系，不能只看人工粘性系数中的一个因子。

\subsection{开关变量的影响}

密度、速度和压力都可以作为人工粘性的开关变量，但表现并不相同。Sod 和 LAX 中，density sensor 与 pressure sensor 比较稳定，velocity sensor 的表现较差，Sod 基准条件下 velocity sensor 甚至失败。第四组超声速膨胀则相反，能稳定完成的组合主要是 velocity sensor。

这说明开关变量的选择和具体波系有关。激波管中密度和压力跳跃能较好标记间断位置；近真空膨胀问题中，密度和压力本身可能非常小，速度变化反而更适合作为稳定化的指示量。

\subsection{网格加密的影响}

开启人工粘性时，网格加密通常可以降低 $L_1$ 误差。例如 Sod 问题在 $\beta=0.5$、CFL=0.5、sensor=$\rho$ 时，密度 $L_1$ 从 $N_x=101$ 的 $1.079602\times10^{-2}$ 降到 $N_x=1001$ 的 $3.227975\times10^{-3}$。LAX 问题也有类似趋势。

但是，网格加密不能从根本上消除 MacCormack 格式处理间断时的色散问题。Sod 问题在关闭人工粘性时，$N_x=101,201,501,1001$ 全部失败。LAX 问题虽然无人工粘性可以完成，但总变差 TV 随网格数增加反而变大，说明间断附近仍有较尖锐的局部结构或振荡。

\section{综合讨论}

七组算例可以分成三类来看。

第一类是 Sod、LAX 和双激波这类含激波问题。MacCormack 格式可以捕捉波系的大体位置，但激波附近会出现局部误差。人工粘性可以减少非物理振荡，但会抹宽间断。Sod 对人工粘性的依赖最明显，关闭人工粘性后无法稳定完成。

第二类是双膨胀和超声速膨胀问题。膨胀问题的困难主要来自低密度、低压力区域。第三组在降低 CFL 后可以稳定计算；第四组接近真空，只有很少参数组合可以完成，是本次七组中最敏感的一组。

第三类是接触间断问题。第五组和第七组都能稳定完成。特别是纯接触间断中，速度和压力误差接近零，主要误差来自密度间断的扩散。这说明不能简单地说某个格式在所有问题上都更好，不同格式对不同波系的优势并不相同。

\section{结论}

本文用带人工粘性的 MacCormack 格式完成了七组一维 Riemann 问题计算。主要结论如下：
\begin{itemize}
    \item MacCormack 格式在光滑区和接触间断问题上表现较好，但在激波和强膨胀区域容易产生色散振荡。
    \item 人工粘性对 Sod 和近真空膨胀问题非常重要，可以避免负压力、负密度和 NaN；但在 LAX、双激波等已经稳定的算例中，人工粘性可能增大整体误差。
    \item $\beta$ 存在折中：太小可能不足以稳定激波，太大又会抹宽间断。本文算例中 $\beta=0.5$ 对多数普通算例较稳，第四组近真空问题需要更强粘性。
    \item CFL 对稳定性和误差的影响不是单调的。虽然人工粘性中有 $\mathrm{CFL}(1-\mathrm{CFL})$，但实际稳定性还受时间步长本身影响。
    \item 开关变量没有固定最优选择。多数激波管问题中密度和压力开关更稳，近真空超声速膨胀中速度开关更容易计算成功。
    \item 单纯加密网格不能完全消除 MacCormack 格式处理间断时的振荡，特别是在关闭人工粘性时，这一点在 Sod 问题中最明显。
\end{itemize}

\begin{thebibliography}{9}
\bibitem{maccormack1969}
R. W. MacCormack.
The effect of viscosity in hypervelocity impact cratering.
AIAA Paper 69-354, 1969.

\bibitem{sod1978}
G. A. Sod.
A survey of several finite difference methods for systems of nonlinear hyperbolic conservation laws.
\textit{Journal of Computational Physics}, 27(1):1--31, 1978.

\bibitem{toro2009}
E. F. Toro.
\textit{Riemann Solvers and Numerical Methods for Fluid Dynamics}.
Springer, 3rd edition, 2009.
\end{thebibliography}

\end{document}
